\textbf{Changes in Rank of Highly Cited Papers and Obsolescence}\\
Although highly cited papers maintain top ranks for a long time owing to the cumulative effect of citations, eventually a sufficient decay in citations is observed, and the papers relinquish their positions to newer, cumulativelygenerated generations of papers. The cumulative citation effect that precedes this turnover is often very long, with the longest observed duration exceeding 50years. In the set of papers for which a clear decay was detected, the characteristics of citationcount changes (the mode from publication to the citation peak and the shape of longterm citation trajectories) were consistent with previously reported findings. Conversely, some papers did not exhibit a marked decay; for these, the interval from publication to the citation peak tended to be longer, and the topranked papers belong to this group.

\textbf{Characteristics of Journals Publishing Highly Cited Papers}\\
When the definition of highly cited paper incorporates longterm citations, it becomes evident that such papers are not necessarily published in journals with high impact factors. The impact factors of journals that host highly cited papers range widely, from 3 to 50. By contrast, journals classified as G2G4 have impact factors clustered around 310, and no journals with particularly high impact factors were found in that group. Moreover, even though journals with higher impact factors do exist, this study did not observe any of them among the venues of the topcited papers.

\textbf{Highly Cited Topics and the Citation Topics of Highly Cited Papers}\\
Highly cited papers tend to focus on a relatively narrow research topic, yet they are cited frequently by papers covering a variety of different topics. When examined separately, the diversity (entropy) of both the research topics and the citing topics of highly cited papers is low; however, the degree of topic switching between the research topic and the citing topic is comparatively high. If such switching influences obsolescence, then papers whose topics attract interest from many research domains may enjoy greater longevity. In line with the journal characteristics mentioned above, journals that publish highly cited papers include those from broad, fundamental or interdisciplinary fields that can draw attention across many areas, whereas journals dedicated to specific applied domains such as medicine or biomedical engineering are absent, supporting this conjecture.

\textbf{LongTerm Trends} \\
The diversity (entropy) of research topics / citing topics shows an upward trend over time in all groups, with highly cited papers consistently exhibiting lower diversity than the other groups. The magnitude of topic switching between research and citing topics remains stable over time across all groups, though highly cited papers display a slightly different pattern. The extent of topic switching between successive cumulative generations is not stable, and no particular trend is discernible. Likewise, the rates of novel appearances of highly cited papers, their publishing journals, and their topicsthat is, the degree of rank turnoverare not stable, and no specific trend is observed.
